\chapter{Introduction}\label{ch:introduction}

\section{A Motivating Scenario}\label{sec:ch1_motivating_scenario}

Imagine a crowded Scottish pub where people have gathered to watch football. Suppose that one infectious person transmits a respiratory virus to several others during the evening. Those people then return to different households and communities. Some later test positive, and samples from their infections are sequenced. This is a hypothetical shared-exposure event; the question for surveillance is whether it could be recognised from the resulting data.

Without records of attendance or contact, the sampled infections may appear unrelated. Their genomes may be very similar, but similar viruses may also be circulating among people who did not attend the gathering. Differences in when people develop symptoms, seek testing, and provide samples further obscure the connection. Genomic similarity and sampling dates could help identify a group worth investigating, while leaving its source and transmission history uncertain.

This scenario motivates the central question of the thesis: \emph{how can genomes and sampling dates be used to assess recent linkage between infections, identify unusual cluster patterns, and guide investigation of possible superspreading?} Answering it requires a clear account of what the genomic relationships represent and how surveillance practices shape the infections that are observed.

% \clearpage
\section{Pathogen Genomics for Public Health Surveillance}\label{sec:ch1_genomic_surveillance}

Genomic surveillance combines pathogen sequences with diagnostic, epidemiological, and contextual information to support public health investigation \citep{WHO2023,Armstrong2019,Black2020}. As pathogens spread, genetic changes create patterns of relatedness among sampled infections. These patterns can help investigate outbreaks, track introductions and geographic spread, and monitor circulating variants \citep{Gardy2017,Struelens2024}. Their interpretation depends on information about when and where samples were collected and how they entered surveillance.

The coronavirus disease 2019 (COVID-19)\nomenclature{COVID-19}{coronavirus disease 2019} pandemic made this relationship between genomic data and surveillance especially important. Large sequencing programmes created opportunities to study severe acute respiratory syndrome coronavirus 2 (SARS-CoV-2)\nomenclature{SARS-CoV-2}{severe acute respiratory syndrome coronavirus 2}, the virus causing COVID-19, at population scale \citep{Brito2022,Carter2022}. Yet sequenced infections were an incomplete sample of the epidemic. Testing access, reasons for testing, sequencing priorities, and the availability of linked information about cases and samples affected which infections became visible. These conditions could change over time and differ between places and population groups.

Superspreading provides a particular motivation for interpreting these data. Transmission can be highly uneven, with a minority of infected individuals causing many more secondary infections than others \citep{LloydSmith2005,Endo2020}. A rapid increase in samples with similar genomes may therefore signal concentrated transmission. However, a large group of similar genomes can also reflect several transmission generations, repeated introductions of a common virus, or intensive sampling of an outbreak. Genomic clusters need to be interpreted alongside the epidemiological and surveillance conditions that produced them.

Population context matters for the same reason. Studies of COVID-19 documented inequalities in infection and health outcomes across demographic, socioeconomic, and geographic groups \citep{Bambra2020,Niedzwiedz2020,Mathur2021,PublicHealthScotland2020}. These findings motivate examining which populations are represented in genomic surveillance and in the clusters inferred from it. Such comparisons describe the sampled record. Establishing individual infection risk, exposure settings, or the causes of population differences requires additional evidence.

This thesis develops a method for assessing pairwise genomic relationships and applies it to Scottish SARS-CoV-2 surveillance. The Scottish analysis proceeds through three stages: describing the sampled background, screening for unusual cluster patterns, and characterising the populations represented in those patterns. Chapter~\ref{ch:literature} reviews the literature and conceptual background supporting this progression.

% \clearpage
\section{Recent Linkage and Genomic Cluster Screening}\label{sec:ch1_recent_linkage}

The focus of this thesis is \emph{recent transmission linkage}, shortened here to \emph{recent linkage}. The term can describe direct transmission between two sampled infections, infections from a shared source, or infections linked by a short chain involving unsampled intermediates. The analyses in this thesis use the base case: direct transmission or co-primary infection, meaning that both sampled infections directly acquired the pathogen from the same source. Section~\ref{sec:ch1_inferential_target} formalises this choice.

Recent linkage is difficult to assess from consensus genomes. A \emph{consensus genome} is a representative sequence that typically records the most common nucleotide at each genomic position in a sample, without capturing the full diversity of lower-frequency variants \citep{Valesano2021}. Early estimates placed the SARS-CoV-2 substitution rate at approximately $10^{-3}$ substitutions per site per year, equivalent to roughly two to three substitutions per genome per month \citep{Duchene2020}. Transmission can therefore occur without a detectable change in the consensus sequence, limiting the ability of genetic similarity alone to distinguish individual transmission links \citep{Braun2021}.

Nearly identical genomes may also come from different transmission chains within a widely circulating group of related viruses (a viral lineage) \citep{Louca2021,GallegoGarcia2022,Fiam2024}. Conversely, recently linked infections need not have identical genomes. Genetic distance therefore needs to be interpreted alongside sampling time, pathogen evolution, and the relationships being considered.

Here, \emph{compatibility} describes how well the observed genetic and temporal differences agree with simulations under specified linkage scenarios and model assumptions. It does not identify who infected whom or provide a probability that transmission occurred. The scores define a network, or \emph{graph}. Each node represents a sampled genome, and each retained connection has a weight equal to the compatibility score for that pair. Clustering groups sampled genomes according to the pattern of these pairwise connections. A cluster can therefore support a hypothesis of recent shared transmission without identifying a single exposure event or a complete transmission chain.

The Scottish analyses repeat this process over overlapping periods of sample collection, termed \emph{surveillance windows}. Clusters are linked across time when they contain some of the same sampled genomes. These links describe continuity in the observed data, without establishing transmission from one cluster to another.

Screening focuses on two patterns. The \emph{local burst} score combines two features: how large a cluster is compared with other clusters sampled under comparable conditions, and the proportion of members absent from directly linked earlier clusters, where those links are observed. Here, `local' refers to the cluster and its immediate links in the data. The \emph{onward burden} score measures additional sampled infections in later linked clusters relative to the starting cluster's size. Clusters flagged by the screen are termed \emph{candidates} for epidemiological investigation; this does not confirm superspreading. The populations represented in these clusters are then examined.

% \clearpage
\section{Transmission Scenarios Included in the Compatibility Score}\label{sec:ch1_inferential_target}

For two sampled infections $i$ and $j$, plausible transmission histories can be represented by a set of scenarios allowing up to $M$ unsampled intermediates (Figure~\ref{fig:transmission_scenarios}):

\[
    \mathcal{S}_M =
    \Bigl\{\, H_{\mathrm{AD}}(m) \mid m \in \{0,\dots,M\} \,\Bigr\}
    \cup
    \Bigl\{\, H_{\mathrm{CA}}(m_i,m_j) \mid m_i,m_j \in \mathbb{Z}_{\ge 0},\;
    m_i+m_j \le M \,\Bigr\}.
\]

Here, $H_{\mathrm{AD}}(m)$ describes transmission along a chain from one sampled infection to the other through $m$ intermediates. The scenario $H_{\mathrm{CA}}(m_i,m_j)$ describes infection from a shared unsampled source, with $m_i$ and $m_j$ intermediates on the two branches. The shared source itself is distinct from these intermediates.

\begin{figure}[htbp]
    \centering
    \includegraphics[width=\textwidth]{transmission_scenarios}
    \caption[Plausible transmission relationships between sequenced infection pairs]{\textbf{Plausible transmission relationships between sequenced infection pairs.} Solid nodes indicate sampled infections; dashed nodes indicate unsampled infections. Symbols labelled $m$, $m_i$, and $m_j$ represent counts of unsampled intermediates; zero means no intermediary. In the common-ancestor scenario, $H_{\mathrm{CA}}(m_i,m_j)$, both infections descend from the shared unsampled source $a$, which is separate from these counts.}\label{fig:transmission_scenarios}
\end{figure}

The selected subset $\mathcal{S}_\star \subseteq \mathcal{S}_M$ contains the transmission scenarios included in the compatibility score. The primary analyses here use $M=0$ and

\[
    \mathcal{S}_\star =
    \bigl\{H_{\mathrm{AD}}(0), H_{\mathrm{CA}}(0,0)\bigr\}.
\]

This includes direct transmission and co-primary infection from a shared unsampled source, with no additional intermediates. Longer chains remain possible explanations, but are not among the scenarios used to calculate the score. A high score indicates that the observations fit the selected scenarios under the model assumptions; it does not exclude alternative histories. Chapter~\ref{ch:epilink} develops and evaluates the corresponding compatibility measure.

% \clearpage
\section{Aims, Research Questions, and Objectives}\label{sec:ch1_aims_questions_objectives}

The thesis has three aims. They connect method development with the screening and epidemiological characterisation of genomic patterns in Scottish surveillance data.

\begin{enumerate}[
    label=\textbf{\arabic*.},
    leftmargin=1.8em,
    labelsep=0.6em,
    topsep=0.5\baselineskip,
    itemsep=0.6\baselineskip,
    parsep=0.15\baselineskip,
    before=\raggedright%
]
    \item {\raggedright\bfseries Develop and evaluate a method for assessing pairwise compatibility with recent linkage.\par}\nopagebreak
    \textit{Question:} How can genetic differences and sampling dates be combined to assess whether two sampled infections are compatible with specified recent-linkage scenarios?

    \item {\raggedright\bfseries Define recent-linkage clusters and screen their patterns over time for possible concentrated transmission.\par}\nopagebreak
    \textit{Question:} Which patterns of cluster size and membership over time are unusual compared with other clusters observed under similar surveillance conditions and warrant investigation for possible superspreading?

    \item {\raggedright\bfseries Characterise the populations and surveillance contexts represented in the inferred graphs, clusters, and screening outcomes.\par}\nopagebreak
    \textit{Question:} How do these patterns vary by demographic, socioeconomic, and geographic attributes, and what questions do those differences raise for epidemiological investigation?
\end{enumerate}

Chapter~\ref{ch:epilink} addresses method development and evaluation. The Scottish application follows seven linked objectives, shown in Figure~\ref{fig:objectives}. The first four establish the surveillance context and construct graphs and clusters from pairwise compatibility scores. The next two link clusters across time and screen for unusual patterns. The final objective examines which population groups are represented in clusters and how diverse those groups are, in relation to the screening outcomes. At the graph level, characterisation asks whether higher compatibility scores tend to connect sampled infections in the same population or residential categories. This helps define the background for interpreting the screening results.

\definecolor{stage1}{HTML}{EAF3FF}
\definecolor{stage2}{HTML}{DDF5F3}
\definecolor{stage3}{HTML}{E3F6E8}
\definecolor{stage4}{HTML}{F1F7D5}
\definecolor{stage5}{HTML}{FFF0C7}
\definecolor{stage6}{HTML}{FFE1CC}
\definecolor{stage7}{HTML}{EDE7F6}

\begin{figure}[p]
    \centering
    \thesisflowbox{stage1}{\textbf{1. Surveillance data and context}\\
        \textit{Objective:} Assemble linked Scottish genomic data and describe the sampled population and surveillance conditions.}
    \thesisflowarrow%
    \thesisflowbox{stage2}{\textbf{2. Time and lineage comparisons}\\
        \textit{Objective:} Organise genomes into surveillance windows and viral lineages.}
    \thesisflowarrow%
    \thesisflowbox{stage3}{\textbf{3. Pairwise compatibility graphs}\\
        \textit{Objective:} Score pairwise relationships and examine whether higher compatibility scores tend to connect samples in the same population or residential categories.}
    \thesisflowarrow%
    \thesisflowbox{stage4}{\textbf{4. Within-window clusters}\\
        \textit{Objective:} Group sampled genomes using their pairwise compatibility relationships.}
    \thesisflowarrow%
    \thesisflowbox{stage5}{\textbf{5. Cluster continuity across time}\\
        \textit{Objective:} Connect clusters across time using sampled genomes shared between windows.}
    \thesisflowarrow%
    \thesisflowbox{stage6}{\textbf{6. Screening for unusual patterns}\\
        \textit{Objective:} Flag clusters with unusual local burst or onward burden scores for investigation.}
    \thesisflowarrow%
    \thesisflowbox{stage7}{\textbf{7. Population characterisation}\\
        \textit{Objective:} Examine which population groups are represented within clusters, how mixed those groups are, and how these profiles relate to the screening results.}

    \caption[Objectives for the Scottish analyses]{\textbf{Objectives for the Scottish analyses.} EpiLink, developed and evaluated in Chapter~\ref{ch:epilink}, supplies the pairwise measure. Objectives 1--4 establish the surveillance background and construct graphs and clusters; Objectives 5--6 identify unusual patterns across time; Objective 7 examines their population profiles. The stages support epidemiological investigation without treating flagged clusters as confirmed superspreading events.}\label{fig:objectives}
\end{figure}

% \clearpage
\section{Thesis Contributions}\label{sec:ch1_contributions}

The thesis makes four connected methodological and empirical contributions:

\begin{enumerate}[
    label=\textbf{\arabic*.},
    leftmargin=1.8em,
    labelsep=0.6em,
    topsep=0.5\baselineskip,
    itemsep=0.6\baselineskip,
    parsep=0.15\baselineskip,
    before=\raggedright%
]
    \item {\raggedright\bfseries A defined measure of recent-linkage compatibility.\par}\nopagebreak
    EpiLink compares genetic and sampling-time differences with simulations under specified transmission scenarios. Its evaluation uses synthetic data with known links and a Boston genomic dataset containing cases linked to documented outbreaks. The method provides scores that express degrees of compatibility between sampled infections without fitting to pairs with known transmission links, while making its assumptions and interpretation explicit.

    \item {\raggedright\bfseries Patterns of genomic compatibility in Scottish surveillance.\par}\nopagebreak
    The Scottish analysis applies EpiLink to 322,366 high-quality genomes across 134 overlapping surveillance windows. Higher compatibility scores are concentrated between samples with matching residential categories: health board, local authority, and urban/rural class. Graph size varies with the numbers of sequenced genomes and recorded positive tests. These findings establish the background against which unusual clusters can be assessed and interpreted.

    \item {\raggedright\bfseries A calibrated screen that assesses local bursts and onward burden separately.\par}\nopagebreak
    Clusters are linked across time through the sampled genomes they share. Calibration assesses how unusual each score is relative to comparable clusters. The local burst and onward burden scores select largely different candidates, providing distinct questions for further investigation.

    \item {\raggedright\bfseries Populations represented in the screened clusters.\par}\nopagebreak
    The analysis examines which population groups are represented in clusters and how diverse those groups are within each cluster. It identifies modest demographic differences between candidate and background clusters. The two screening scores also show different associations with the mix of residential health boards within clusters, after accounting for cluster size and the population sampled at the time. These findings give the signals an epidemiological context while retaining the distinction between sampled associations and explanations of transmission.
\end{enumerate}

Taking these contributions together, this thesis defines an advancement in formalised, \emph{process-based}, genomic surveillance. This approach interprets genetic differences and sampling dates using explicit models of transmission and mutation. It assesses the resulting clusters alongside the conditions that determined which infections were tested and sequenced. This framework connects an assessment of pairwise relationships to the screening and characterisation of patterns that can be investigated using epidemiological records.

% \clearpage
\section{Thesis Structure}\label{sec:ch1_structure}

The thesis moves from the problem and its conceptual foundations to method development, three empirical analyses, and an overall synthesis.

\begin{description}[
    style=multiline,
    leftmargin=6.4em,
    labelwidth=5.8em,
    labelsep=0.6em,
    align=left,
    font=\bfseries,
    topsep=0.5\baselineskip,
    itemsep=0.6\baselineskip,
    parsep=0.15\baselineskip,
    before=\raggedright%
]
    \item[Chapter~\ref{ch:introduction}] Introduces the motivating problem, defines the recent-linkage target, and sets out the aims, questions, and contributions.

    \item[Chapter~\ref{ch:literature}] Reviews what genomic surveillance observes and compares approaches to linkage and clustering, establishing the basis for screening cluster patterns over time and population characterisation.

    \item[Chapter~\ref{ch:epilink}] Develops and evaluates EpiLink. It tests how well the method distinguishes pairs meeting the recent-linkage definition from other pairs, and compares inferred clusters with reference groups in simulations and documented outbreak exposures in Boston. It also examines how results respond to modelling assumptions and additional data.

    \item[Chapter~\ref{ch:genomic_networks}] Begins the Scottish empirical analysis. It describes the surveillance record, constructs compatibility graphs within windows and lineages, and examines their demographic, geographic, and socioeconomic structure.

    \item[Chapter~\ref{ch:cluster_transition_screening}] Links clusters across time through shared sampled genomes. It introduces the local burst and onward burden scores and calibrates them against comparable clusters to identify candidates for investigation.

    \item[Chapter~\ref{ch:candidate_characterisation}] Examines population composition and diversity in relation to candidate status and screening scores. Bayesian hierarchical models assess demographic, geographic, and surveillance associations, allowing model baselines to differ between policy periods and viral clades (broader groups of related viral lineages).

    \item[Chapter~\ref{ch:discussion}] Brings the findings together around the three research questions. It considers the combined epidemiological and methodological contribution, the limits of interpretation, and priorities for external validation and practical evaluation.

    \item[Appendices] Appendix~\ref{app:epilink} gives EpiLink's derivation and supporting evaluations. Appendix~\ref{app:genomic_networks} covers data processing, graph-construction checks, and supplementary analyses. Appendix~\ref{app:cluster_transition_screening} gives technical details and robustness checks for screening cluster patterns over time. Appendix~\ref{app:candidate_characterisation} contains Bayesian model specifications, diagnostics, and extended characterisation results.
\end{description}
